%% ch04 --- Motion: rules of change and the map of every state
%%
%% Stub, generated by tools/gen_stubs.py from tools/chapters.json.
%% The brief below is the contract for this chapter: what it must argue, what
%% it must buy the reader, and what it must NOT take from its neighbours.
%% Writing the chapter means deleting the stub block below (the one that opens
%% on the next \chapter line) and replacing it with the chapter text -- after
%% which this file is never regenerated.
%% If the brief turns out to be wrong once the code has been run, change it in
%% the manifest and record the contradiction in CLAUDE.md under *Resolved
%% questions*.

\chapter{Motion: rules of change and the map of every state}
\label{ch:motion}

\chapterstub{%
Argument: when time flows rather than ticks, the rule tells you the VELOCITY
of the state rather than its next value \dash{} a differential equation,
read as a rule of change and never solved by hand here. Euler's method turns
it back into iteration with a small step; the spring and the pendulum are
the worked systems. The state of a pendulum is two numbers, angle and
angular velocity, so its whole future is a curve through a plane of every
possible state: the phase space, the one picture the rest of the book is
drawn in. A frictionless pendulum draws closed loops (periodic, perfectly
predictable); a damped one spirals into a fixed point (an attractor). Energy
is the check: Euler's method visibly gains energy and spirals OUTWARD, a
better method (Runge-Kutta, RK4, given as a recipe) holds it \dash{} so the
reader learns early that a simulation is not the system. Newton's two bodies
trace ellipses forever: the clockwork at its best. Plant the theorem of
Poincaré and Bendixson in one sentence: a smooth system with a two-number
state cannot be chaotic, so chaos in a flow needs at least three numbers
\dash{} Chapter 9 cashes it. Payoff: the reader can turn a rule of change
into a trajectory in code and read a phase portrait. Traps to elicit: a
smaller step always makes a simulation right (it makes it right for longer);
a pendulum is a simple system so its motion must be simple (true for one
pendulum, false for two, Chapter 13). Listings: Euler and RK4 on a spring;
the pendulum; energy drift. Labs: implement one RK4 step; measure energy
drift. Figures: phase portrait of a pendulum; Euler spiralling out against
RK4.

\medskip
\textbf{Word budget:} 3500.
}
